% Editable visual-geometry calculation note for issue #15901.
% This is a stylized, source-informed appearance model, not an anatomical fit.
\documentclass[11pt]{article}
\usepackage[margin=1in]{geometry}
\usepackage{amsmath,amssymb}
\usepackage{booktabs}
\usepackage{hyperref}
\title{Tousled-Wave Hair on the Closed CC0 Scalp}
\author{UpstreamDrift appearance research}
\date{2026-11-28}
\begin{document}
\maketitle

\section{Purpose and evidence}
This note defines the visual-only \texttt{tousled\_wave} style for the CC0
head's existing scalp shell. The source-inspired target is the full-volume,
short wavy hair visible in the contemporary PGA TOUR portraits of Mac O'Grady
and John Schlee. The 1984 PGA TOUR Guide describes O'Grady's profile and
reproduces a grayscale headshot (printed page 124); the 1973 PGA TOUR Guide
shows Schlee's grayscale portrait and calls him curly-haired (printed page
58). These are qualitative references. They do not provide calibrated head
geometry, color, or strand-level measurements.

\begin{itemize}
  \item O'Grady: \url{https://pgatourmedia.pgatourhq.com/static-assets/page/header/files/1984_pgatour.pdf}
  \item Schlee: \url{https://pgatourmedia.pgatourhq.com/static-assets/page/header/files/1973_pgatour.pdf}
\end{itemize}
The source images are retained as local research material only; neither image
is copied into this repository. Grayscale tone does not establish natural hair
color, and the headshots do not support a body-build inference.

\section{Coordinates and inputs}
All length inputs are in head-length units, with the existing CC0 head's neck
plane at $z=0$ and vertex at $z=1$. In the canonical frame, $x$ points toward
the face, $y$ points to the subject's left, and $z$ points upward. The existing
closed scalp mesh supplies the vertices $\mathbf{p}_i$ and outward normals
$\mathbf{n}_i$. A style adds only a bounded shell around this surface.

Let $s\in\{-1,0,+1\}$ encode the hair part. The phase of the existing comb
ridges is
\begin{align}
 u(x,y) &= \frac{y-s\,q(x-x_c)}{\lambda_c},\\
 \theta(x) &= \frac{2\pi(x-x_c)}{\lambda_w},\\
 b(y) &= 0.18\sin\left(\frac{2\pi y}{\lambda_c}\right),\\
 \phi(x,y) &= 2\pi\left[u(x,y)+a_w\left(\sin(\theta+b)+0.20\sin(2\theta+0.7-b)\right)\right].
\end{align}
where $q$ controls ridge sweep and $b(y)$ varies ridge bending across the scalp. The style sets $a_w=0.42$ cycles and
$\lambda_w=0.26$ head lengths. The wave term bends the existing ridges across
the crown; it does not add disconnected locks, vertices, or a second shell.
For every other hairstyle, $a_w=0$ and the phase-warp branch is skipped, which
preserves its previous vertex coordinates exactly.

The added thickness field is
\begin{equation}
 t_w(x,y,z)=A_c\,M_{\mathrm{crown}}(z)\cos(\phi(x,y)),
\end{equation}
with crown mask $M_{\mathrm{crown}}$ already used by the existing haircut
builder. The full thickness is the current base fade, fringe, and parting
field plus $t_w$. It is multiplied by the existing smooth hairline edge mask,
so added volume returns to zero at the cut boundary. The style-specific inputs
are:
\begin{center}
\begin{tabular}{lll}
\toprule
Input & Value & Meaning \\
\midrule
$A_c$ & 0.015 & comb-ridge amplitude (head lengths) \\
$\lambda_c$ & 0.10 & ridge wavelength across the crown (head lengths) \\
$a_w$ & 0.42 & sinusoidal phase warp (cycles) \\
$\lambda_w$ & 0.26 & bend wavelength along the crown (head lengths) \\
$t_{\mathrm{top}}$ & 0.065 & top shell thickness (head lengths) \\
$t_{\mathrm{side}}$ & 0.018 & side shell thickness (head lengths) \\
$t_{\mathrm{fringe}}$ & 0.046 & soft front-hairline lift (head lengths) \\
front-edge wave & 0.006 & hairline-height amplitude (head lengths) \\
\bottomrule
\end{tabular}
\end{center}
The base hairline control points retain a modest forehead and longer nape;
parting depth and heavy-side gain are zero. The hairline receives a bounded
0.006 head-length sinusoidal variation over the front arc; it is subtle and
does not change the closed rim. \texttt{tousled\_wave} remains nearly symmetric. Its top remains below $z=1.09$ in the unit-head geometry.

\section{Closed-shell construction and invariants}
For each scalp point, the existing mesh builder evaluates the nonnegative
total thickness and offsets the outer surface along $\mathbf{n}_i$. It keeps
the hidden inner surface at the configured skin clearance and zips the shared
hairline rim to close the shell. The wave changes only the scalar thickness;
it does not change the scalp cut, inner surface, face winding, or rim index
welding. The result must remain a finite, outward-wound, closed manifold with
positive signed volume.

When a cap is requested, the existing cap-specific path still builds the
legacy short crop and uses its cap-clearance geometry. The new style does not
alter that fallback. Appearance has no effect on the native body, joint,
contact, or dynamics arrays; the native MuJoCo regression compares those
arrays and forward-dynamics quantities across all named styles.

\section{Limits and reproduction}
The portrait references motivate the broad silhouette (fuller crown and
wavy/tousled ridges) only. This is not a reconstruction of any named player's
individual hair or a measured historical hairstyle. Colors remain an artistic
choice. To reproduce the geometry checks from the repository root, run:
\begin{verbatim}
python -m pytest tests/unit/model_appearance/test_hair_styles.py -q
\end{verbatim}
The native-model equality test is visual-only qualification: it verifies that
hair changes do not alter modeled dynamics. It is not a motion-fit or physical
validation of either golfer.
\end{document}
